\section{Evaluation}
\label{sec:eval}

\textbf{Setup.}
The corpus spans ten domains and bases, with 138 documents, 879 parts,
and 23.9 MB of source text. We ran 72 in-domain gold questions and four
out-of-domain probes intended to exercise evidence rejection; two
cross-domain placeholders were not run.
Each arm has one trace per question. \emph{A0} is the full retrieval
ladder with description-based routing and the Laya gate; \emph{A1}
mechanically reads all 879 parts without answering; \emph{A1b} uses
global vector recall plus the gate; \emph{A2} removes the gate; and
\emph{A3} removes catalog descriptions. A1b and A2 used an unscoped
global index whose top results include the A3 twin and other external
collections. We report them as exploratory, not controlled baselines.

The primary metric is exact containment of the gold string in packed
text, a retrieval proxy rather than answer correctness. Table
~\ref{tab:ablation} reports gold-only rates with Wilson 95\% intervals,
mean read volume, and evidence-gate rejections. We use exact, qid-paired McNemar
tests only for same-scope comparisons, without multiple-test correction.
The 380 saved traces record tool calls, routing, evidence bytes, gate
outcomes, and elapsed tool time.

\textbf{Containment and cost.}
A1 contained all 72 gold strings; A0 contained 37/72 (51.4\%). A1 read
24.3 MB per question, about 774 times A0's 31.4 KB, and took 253.7 s
of tool time versus 16.3 s. A1 is a mechanical full-corpus reference,
not an answering system. The unscoped A1b and A2 runs contained gold
strings in 42/72 (58.3\%) and 35/72 (48.6\%); their scope prevents a
fair comparison with A0. A3 contained 33/72 (45.8\%). No controlled
end-to-end result establishes retrieval superiority.

The LLM-free BM25 floor contained the gold string in 24/40 questions;
A0 contained it in 23/40 on the same subset. This is descriptive
containment only: BM25 uses non-PDF 2,000-character windows, lowercase
word/CJK tokenization, Okapi $k_1{=}1.2$, $b{=}0.75$, average length
400 tokens, and top-10 packing; it has no evidence gate. Dense and RRF
floors are omitted because their global index returned external
collections.

\textbf{Routing and evidence-gate outcomes.}
Description-based routing was correct for 72/72 questions, versus
37/72 (51\%) for name-only A3. On matched questions, descriptions
corrected 35 routes with no reverse case (exact McNemar $p<.0001$).
Legacy cards scored 34/66 (52\%) in a separate run. On four
out-of-domain probes, A0's evidence gate returned no survivors on 2/4;
A1b and A3 did so on 0/4. A1 has no gate or answer decision, and A2
disables the gate. These are retrieval-stage outcomes, not measured
assistant abstentions; the ablation does not generate answers.

Six hinted native-agent runs completed with correct citations and key
facts; separate verbatim audits checked 23/23 citations, 18/18 quotes,
and 14/14 quotes. These are illustrative runs, not unhinted routing
estimates.

\textbf{Limitations.}
This is a single run on one domain-skewed corpus. Containment does not
measure answer completeness, semantic correctness, or citation
correctness; only four out-of-domain probes were used. A source-PDF audit
found the exact strings for all 14 PDF-backed gold rows after whitespace
normalization; their parsed and indexed parts were not independently
audited. The corpus is domain-skewed, and 52/72 questions were
automatically expanded. A clean, evaluation-scoped dense baseline with
the same evidence gate remains to be run. The results support a
description-sensitive routing effect, not general retrieval superiority.
